%!TEX program = uplatex
\RequirePackage{plautopatch}
\documentclass[uplatex,dvipdfmx]{jlreq}
\usepackage{amsmath,amssymb,amscd,amsthm}

\pagestyle{empty}

\begin{document}

%======================================================================
% 表紙
%======================================================================
\begin{center}
  \vspace*{20pt}
  {\Huge\bfseries 流体力学}

  \vspace{1pc}
  {\Large\bfseries ぐるぐるぎりぎりの流体数学}

  \vspace{1pc}
  {\Large 概要}
\end{center}

\vfill

\newpage

%======================================================================
% 講演１：木村芳文「流体数学のたのしみ」
%======================================================================
\begin{center}
  {\Large\bfseries 流体数学のたのしみ}
\end{center}

\bigskip
\hspace*{\fill}{\large\bfseries 木村　芳文（名古屋大・多元数理）}

\bigskip

流体力学は歴史的に見て数々の問題意識を数学に対して提案してきたと思う。
古くは今井先生からお話があるように等角写像の問題を初め、新しくは可積分・
非可積分性のキーワードとしてのソリトン、カオスといった概念も元々は流体における
現象のモデル化から生じてきたケースが多い。
この尽きることの無い流体現象の多様性は流体の持つ非線形性に起因していると言ってよいだろう。

流体の非線形性の本質は流体の持つエネルギーが異なる長さスケール間を移動するという
いわゆるエネルギーカスケードにある。
（これは数学的には基礎方程式をフーリエ変換すると非線形項は convolution を与え、
全波数からの寄与を受けることに対応する。）
このエネルギーカスケードがあるからこそ、流体の研究は難しくかつ面白い。

本講演では流体数学への導入としてこのエネルギーカスケードの問題についていくつか
考えてみたい。

なお、副題の「ぐるぐるぎりぎりの流体力学」について誤解を招かないように補足しておこう。
「ぐるぐる」というのは流体の非線形性のひとつの柱である渦を意味しているが、
まあこれはわかりやすい。
問題は「ぎりぎり」という言葉であろう。
多くの非線形現象はぎりぎりの状況において、あるいは別の言葉でいえば critical な状況において
非常に美しい数学的・特に幾何学的描写を可能にすることがままある。
流体力学でも同様の局面が現れることは想像に難く無い。
そういった意味合いを含めての言葉である。
注意して聞いて頂ければ今回の講師の皆さんはいずれの方も何らかの意味で、
「ぎりぎり」の状況を扱っておられることに気がつかれるであろう。

\newpage

%======================================================================
% 講演２：今井功「流体力学における等角写像の応用とその限界 --- 私的回想」
%======================================================================
\begin{center}
  {\Large\bfseries 流体力学における等角写像の応用とその限界\\
  \normalsize --- 私的回想}
\end{center}

\bigskip
\hspace*{\fill}{\large\bfseries 今井　功}

\bigskip

物理中期の「物理学演習」で坂井卓三先生が「Helmholtz の噴流」の問題をとり扱われたとき、
その巧妙さに驚嘆したのが、私の「等角写像」とのつきあいのはじまりである。
卒業後、助手として友近晋先生から最初に与えられた問題が「平板翼の海面効果」であった。
これによって「等角写像」は私にとって終生の友となった。しかし、
「流体力学」という視点がつねに脳裏にあるため、
関心の点は、多分、数学者の方々とはいささか（あるいは大いに）異なるのではないかと思う。
私にとって最大の関心事は、

\medskip
\noindent
(i)\quad 任意の領域 $D, D'$ をたがいに等角写像する正則関数を求めること、\\
(ii)\quad 任意の領域 $D$ の境界 $C$ 上で実数部あるいは虚数部を与えて、
$D$ で正則な関数を求めること、
\medskip

\noindent
であった。

(i) については在来の方法が、いわば「天下り」的であるのに不満を覚えていたからである。
これに対して、「流体力学的発想法の定理」なるものを案出した。

\bigskip

\begin{center}
\begin{minipage}[t]{11cm}
  \textbf{[定理]}\quad
  $z$ 平面、$\zeta$ 平面に領域 $D, D'$ が与えられている。
  $f = f(z)$、$F = F(\zeta)$ はそれぞれ $D, D'$ でおこる渦無しの流れを表す
  複素速度ポテンシャルとする。このとき、$f(z) = F(\zeta)$ は
  $D$ を $D'$ に等角写像する正則関数を与える。
\end{minipage}
\end{center}

\bigskip

拙著 [1] はこの定理の応用を例示したものである。実は「等角写像の事典」の類いを
かねがね意図していたのであるが、この定理によって「事典」の効用に疑問を感じるようになった。

なお、(ii) その他の私的回想については [3] をご覧いただければ幸いである。

\vspace{6pt}

\begin{center}
  \textsc{参考文献}
\end{center}
\begin{enumerate}
  \item[{[1]}] 今井　功：『等角写像とその応用』（岩波、1979, 1998）.
  \item[{[2]}] H.~Kober: \textit{Dictionary of Conformal Representations},
               Dover Publications, 1962.
  \item[{[3]}] 今井　功：『複素解析と流体力学』（日本評論社、1989, 1996）.
\end{enumerate}

\newpage

%======================================================================
% 講演３：今井功「数学と物理の交流」
%======================================================================
\begin{center}
  {\Large\bfseries 数学と物理の交流}
\end{center}

\bigskip
\hspace*{\fill}{\large\bfseries 今井　功}

\bigskip

私は数学のアマチュア（愛好者）である。しかし、大学で物理学科に入り、
卒業後も物理屋として過してきた。
その間、数学には大いにお世話になった。しかし、今思うと「道具」としての数学よりも、
むしろその「考え方」に惹かれたのではないかと感じる。

私は、自然現象の探究は
\[
  \text{実在現象}
  \quad \rightleftharpoons \quad
  \text{物理モデル}
  \quad \rightleftharpoons \quad
  \text{数学モデル}
\]
のように進むものと考えている。この観点で「数学」と「物理」の交流を考えてみたい。

\newpage

%======================================================================
% 講演４：宮川鉄朗「非圧縮流体の解析学」
%======================================================================
\begin{center}
  {\Large\bfseries 非圧縮流体の解析学\\
  \normalsize --- 非粘性流体中の渦運動と粘性流体の挙動と ---}
\end{center}

\bigskip
\hspace*{\fill}{\large\bfseries 宮川　鉄朗（神戸大・理）}

\bigskip

流体力学の数学解析から標記の二つの話題を取り上げ、現状を紹介する。

\begin{description}
  \item[(1)] 非粘性流体内の渦の運動については、
             ２次元流の中の vortex patches や vortex sheets の時間発展の追跡について、
             知られた結果を紹介する。
  \item[(2)] 粘性流体の運動については、３次元流の運動方程式の解の存在問題の周辺に
             未解決問題があり、それを克服する努力が今も進行中である。
             現在までに得られた結果の基本的なものを紹介する。
             次いで解の長時間漸近挙動についての最近の結果を紹介する。
\end{description}

\newpage

%======================================================================
% 講演５：吉田善章「Beltrami fields in Fluids and Plasmas」
%======================================================================
\begin{center}
  {\Large\bfseries Beltrami Fields in Fluids and Plasmas}
\end{center}

\bigskip
\hspace*{\fill}{\large\bfseries 吉田　善章（東大・新領域創成科学研究所）}

\bigskip

The Beltrami fields, eigenfunctions of the curl operator,
represent essential characteristics of twisted, spiral, chiral or helical
structures in various vector fields.
Amongst diverse applications of the theory of Beltrami fields,
we discuss the self-organized states and their topological properties.

\end{document}
